\documentclass[11pt]{article}
\usepackage{docstyle}
\renewcommand\sheetname{Light \textperiodcentered{} Handout}

\begin{document}
\sheethead{Light: Unit Review}{Year 9 Science \textperiodcentered{} end-of-year exam}{Handout}

\begin{keybox}[title=The one idea]
Light travels in \textbf{straight lines} from a source. When it meets something it is \textbf{reflected}, \textbf{absorbed}, \textbf{transmitted} or \textbf{refracted}.
Every question in this unit asks two things: \emph{what happens to the light}, and \emph{where does it go next?} Draw the ray, then explain it.
\end{keybox}

\section{How we see}\label{sec:see}
\begin{minipage}[t]{0.55\textwidth}
\vspace{0pt}
A \textbf{light source} gives out (emits) its own light: the Sun, a candle flame, a light bulb, a glow-worm, a phone screen.
A \textbf{light reflector} only reflects light that falls on it: the Moon, a planet, a mirror, this page.

We see an object when \textbf{light from it enters our eyes}. For a non-luminous object the light comes from a source, is reflected by the object, and travels to the eye.

\textbf{Write it in this order} (\emph{Why can we see the Moon?})
\begin{enumerate}
\item The Sun \textbf{emits} light.
\item The light travels in straight lines to the Moon.
\item The Moon \textbf{reflects} some of it.
\item The reflected light travels to Earth and \textbf{enters our eyes}.
\end{enumerate}
The Moon looks dimmer in daytime because the sky is already bright with scattered sunlight; it is \textbf{not} a light source.
\end{minipage}\hfill
\begin{minipage}[t]{0.42\textwidth}
\vspace{2pt}\centering
\FigHowWeSee[7.2mm]
\end{minipage}

\section{Transparent, translucent, opaque}\label{sec:materials}
\begin{center}\FigMaterials[6.2mm]\end{center}
\begin{tabularx}{\textwidth}{@{}l X l@{}}
\toprule
\textbf{Word} & \textbf{What the light does} & \textbf{Examples}\\
\midrule
\tTransparent & almost all light is \textbf{transmitted} straight through; objects behind are seen \textbf{clearly} & window glass, clear water, air\\
\tTranslucent & light passes through but is \textbf{scattered}; objects behind look \textbf{blurred} & baking paper, frosted glass, milk bottle\\
\tOpaque & \textbf{no light} passes through; it is absorbed or reflected & wood, wall, cardboard, jumper\\
\bottomrule
\end{tabularx}

\section{Properties of light}\label{sec:props}
\begin{enumerate}
\item Light travels in \textbf{straight lines}. Evidence: sharp shadows, a laser beam in smoke, a candle seen only through a \emph{straight} tube.
\item Light travels \textbf{very fast}: $300\,000\ \mathrm{km/s}$ in air or a vacuum.
\item Light can travel through a \textbf{vacuum} (empty space), e.g.\ from the Sun to Earth. Sound cannot: it needs particles.
\end{enumerate}
Light slows down in materials (e.g.\ about $200\,000\ \mathrm{km/s}$ in glass, $125\,000\ \mathrm{km/s}$ in diamond) but is still very fast. This change of speed causes refraction (Section~\ref{sec:refraction}).

\textbf{Worked example.} The circumference of the Earth is $40\,000\ \mathrm{km}$. How many times could light go around the Earth in one second?\par
Step 1: in one second light travels $300\,000\ \mathrm{km}$.\par
Step 2: number of times $= \dfrac{300\,000\ \mathrm{km}}{40\,000\ \mathrm{km}} = \mathbf{7.5}$ times.\par
Check: $7.5 \times 40\,000 = 300\,000$. \checkmark

\textbf{Why we see lightning before we hear thunder:} both happen together, but light ($300\,000\ \mathrm{km/s}$) travels about a million times faster than sound (about $0.34\ \mathrm{km/s}$), so the light reaches us first.

\section{Shadows}\label{sec:shadows}
A \textbf{shadow} forms where an opaque object \textbf{blocks} light. Because light travels in straight lines, the shadow has a similar \textbf{shape} to the object.
\begin{center}\FigShadow[5mm]{0}\hspace{8mm}\FigShadow[5mm]{1}\end{center}
\begin{itemize}
\item \textbf{\tUmbra}: the dark, sharp part where \textbf{all} the light is blocked.
\item \textbf{\tPenumbra}: the lighter, fuzzy edge where only \textbf{some} of the light is blocked. It forms only when the source is \textbf{large} compared with the object.
\end{itemize}
\begin{tabularx}{\textwidth}{@{}X X X@{}}
\toprule
\textbf{Situation} & \textbf{What you see} & \textbf{Why}\\
\midrule
small (point) source & sharp shadow, umbra only & every ray from the source is either blocked or not\\
large source & umbra with a penumbra edge & near the edge, part of the source is still visible\\
object moved closer to the source & shadow gets \textbf{bigger} & the edge rays spread out more before the screen\\
\bottomrule
\end{tabularx}
\textbf{Drawing a shadow:} use a ruler; draw straight rays from the \emph{edges} of the source past the \emph{edges} of the object to the screen; add arrows.

\section{Colour}\label{sec:colour}
\subsection*{White light and the spectrum}
White light is a \textbf{mixture} of the colours of the visible spectrum:
\textbf{red, orange, yellow, green, blue, indigo, violet} (\emph{Richard Of York Gave Battle In Vain}). A prism or a diffraction grating splits white light into these colours (Section~\ref{sec:dispersion}).

\subsection*{Primary colours of light}
\begin{minipage}[t]{0.5\textwidth}
\vspace{0pt}
The eye has three types of \textbf{cone} cell, sensitive to \textbf{red}, \textbf{green} and \textbf{blue} light. These are the \textbf{primary colours of light}; they cannot be made by mixing other colours.
Mixing two primary colours gives a \textbf{secondary colour}; all three give \textbf{white}. Phone and TV screens use tiny red, green and blue \textbf{pixels}.

\textbf{Careful:} red, yellow, blue are primary colours of \emph{paint} (pigment). Paint mixing is a different process.
\end{minipage}\hfill
\begin{minipage}[t]{0.47\textwidth}
\vspace{0pt}\centering
\FigRGB[5mm]
\end{minipage}

\subsection*{Coloured objects}
An object \textbf{reflects} its own colour and \textbf{absorbs} all the others. A white object reflects all colours; a black object absorbs all colours.
If the light shining on it contains \textbf{none} of the object's colour, nothing is reflected and the object looks \textbf{black}.
\begin{center}\FigObject[5mm]\end{center}

\textbf{Write it in this order} (\emph{Why does a red apple look black in blue light?})
\begin{enumerate}
\item The apple is red, so it \textbf{reflects only red} light and \textbf{absorbs} all other colours.
\item Blue light contains \textbf{no red} light.
\item The apple \textbf{absorbs} the blue light, so \textbf{no light is reflected} into our eyes.
\item With no light from it reaching our eyes, the apple looks \textbf{black}.
\end{enumerate}

\subsection*{Filters}
A filter \textbf{transmits} (lets through) only its own colour and \textbf{absorbs} the rest. A red filter in white light lets out red light only. A red filter followed by a green filter lets out \textbf{no light}.
\begin{center}\FigFilter[5mm]\end{center}
\begin{tabularx}{\textwidth}{@{}l l X@{}}
\toprule
\textbf{Object colour} & \textbf{Light shone on it} & \textbf{Object appears}\\
\midrule
white & red & \textbf{red} (reflects the red it receives)\\
red & red & \textbf{red}\\
red & green or blue & \textbf{black} (nothing to reflect)\\
blue & white & \textbf{blue}\\
green & magenta (red + blue) & \textbf{black}\\
yellow & red & \textbf{red} (yellow objects reflect red and green)\\
\bottomrule
\end{tabularx}

\section{Reflection}\label{sec:reflection}
\begin{minipage}[t]{0.52\textwidth}
\vspace{0pt}
\begin{itemize}
\item \textbf{\tNormal}: a line drawn at \textbf{90\textdegree} to the mirror where the ray hits it (dashed).
\item \textbf{\tAoI} ($i$): between the incident ray and the \textbf{normal}.
\item \textbf{\tAoR} ($r$): between the reflected ray and the \textbf{normal}.
\end{itemize}
\textbf{Law of reflection:} the \tAoI{} \textbf{equals} the \tAoR: $i = r$.

\textbf{Worked example.} A ray hits a mirror at 25\textdegree{} to the mirror surface. Find $i$ and $r$.\par
Step 1: angles are measured from the normal, which is at 90\textdegree{} to the mirror.\par
Step 2: $i = 90^\circ - 25^\circ = \mathbf{65^\circ}$, so $r = \mathbf{65^\circ}$.
\end{minipage}\hfill
\begin{minipage}[t]{0.45\textwidth}
\vspace{4pt}\centering
\FigReflect[5.6mm]
\end{minipage}

\subsection*{The image in a plane mirror}
\begin{minipage}[t]{0.5\textwidth}
\vspace{0pt}
The image in a plane (flat) mirror is:
\begin{enumerate}
\item the \textbf{same size} as the object;
\item the \textbf{same distance behind} the mirror as the object is in front;
\item \textbf{upright} (the right way up);
\item \textbf{\tLatInv} (left and right swapped);
\item \textbf{\tVirtual}: the light only \emph{appears} to come from it; it cannot be formed on a screen.
\end{enumerate}
\textbf{To locate an image:} (1) measure the object's distance from the mirror and mark the image the same distance behind, on a line at 90\textdegree{} to the mirror; (2) draw two rays from the object to the mirror; (3) draw the reflected rays so that, extended backwards as \textbf{dotted lines}, they meet at the image.
\end{minipage}\hfill
\begin{minipage}[t]{0.46\textwidth}
\vspace{2pt}\centering
\FigPlaneImage[5mm]
\end{minipage}

\begin{tabularx}{\textwidth}{@{}l X X@{}}
\toprule
& \textbf{Real image} & \textbf{Virtual image}\\
\midrule
Meaning & light rays \textbf{actually meet} there & light rays only \textbf{appear} to come from there\\
On a screen? & \textbf{yes} & \textbf{no}\\
Examples & projector, camera, image on the retina & plane mirror, magnifying glass\\
\bottomrule
\end{tabularx}

\medskip
\begin{minipage}[t]{0.6\textwidth}
\vspace{0pt}
\textbf{Lateral inversion} is why AMBULANCE is printed back-to-front on the front of the van: a driver ahead sees it the right way round in the rear-view mirror.
\begin{center}\FigLatInv\end{center}
\textbf{Uses of plane mirrors:} a \textbf{periscope} has two parallel mirrors at \textbf{45\textdegree}; light turns through 90\textdegree{} at each mirror, so you can see over or around obstacles. Kaleidoscopes, dentist's mirrors and rear-view mirrors also use reflection.
\end{minipage}\hfill
\begin{minipage}[t]{0.35\textwidth}
\vspace{0pt}\centering
\FigPeriscope[5mm]{1}
\end{minipage}

\section{Refraction}\label{sec:refraction}
\textbf{Refraction} is the \textbf{bending} of light as it passes from one medium into another, because its \textbf{speed changes}. A medium in which light is slower is more \textbf{optically dense}.
\begin{center}\FigRefractRules[5mm]\end{center}
\begin{tabularx}{\textwidth}{@{}X X X@{}}
\toprule
\textbf{Into a slower medium} (air $\to$ glass or water) & \textbf{Into a faster medium} (glass or water $\to$ air) & \textbf{Along the normal}\\
\midrule
slows down, bends \textbf{towards} the normal & speeds up, bends \textbf{away from} the normal & speed changes but \textbf{no bending}\\
\bottomrule
\end{tabularx}
Memory aid from the notes: \emph{Strawberries Taste Nice}: \textbf{S}lower $\to$ \textbf{T}owards the \textbf{N}ormal.

\begin{minipage}[t]{0.49\textwidth}
\vspace{0pt}\centering
\FigBlock[5mm]{0}\par\medskip
\raggedright\small\textbf{Rectangular block:} the ray bends towards the normal going in and away coming out. It leaves \textbf{parallel} to the incident ray but shifted sideways.
\end{minipage}\hfill
\begin{minipage}[t]{0.49\textwidth}
\vspace{0pt}\centering
\FigSemi[5mm]\par\medskip
\raggedright\small\textbf{Semicircular block:} a ray along a radius meets the curved face along the normal, so it does not bend there; it bends only at the flat face.
\end{minipage}

\subsection*{Refraction in everyday life: things under water look closer}
\begin{minipage}[t]{0.5\textwidth}
\vspace{0pt}
\textbf{Write it in this order} (\emph{Why does a fish look closer to the surface than it is?})
\begin{enumerate}
\item Light from the fish travels from water into air.
\item It \textbf{speeds up}, so it bends \textbf{away from the normal} at the surface.
\item Our brain assumes light travels in \textbf{straight lines}, so it traces the rays straight back (dotted lines).
\item The rays appear to come from a point \textbf{higher} than the fish: we see a \textbf{virtual image} closer to the surface.
\end{enumerate}
So a pool looks \textbf{shallower} than it is, a straw looks bent at the water surface, and a spear fisher must aim \textbf{below} where the fish appears.
\end{minipage}\hfill
\begin{minipage}[t]{0.46\textwidth}
\vspace{0pt}\centering
\FigFish[7mm]
\end{minipage}

\section{Dispersion}\label{sec:dispersion}
\begin{minipage}[t]{0.52\textwidth}
\vspace{0pt}
\textbf{\tDispersion} is the \textbf{splitting of white light} into its colours.
In a prism, each colour travels at a \textbf{different speed}, so each is \textbf{refracted by a different amount}.
\textbf{Violet} bends the \textbf{most}; \textbf{red} bends the \textbf{least}.
A \textbf{rainbow} forms when sunlight is refracted and dispersed by raindrops.

\emph{The spread of colours in the diagram is exaggerated so it can be seen.}
\end{minipage}\hfill
\begin{minipage}[t]{0.45\textwidth}
\vspace{0pt}\centering
\FigPrism[5.4mm]
\end{minipage}

\section{The eye}\label{sec:eye}
\begin{minipage}[t]{0.5\textwidth}
\vspace{0pt}\centering
\FigEye[6.2mm]{1}
\end{minipage}\hfill
\begin{minipage}[t]{0.48\textwidth}
\vspace{0pt}
\small
\begin{tabularx}{\linewidth}{@{}l X@{}}
\toprule
\textbf{Part} & \textbf{Function}\\
\midrule
\tCornea & clear protective front; \textbf{refracts} (bends) the light most\\
\tIris & coloured \textbf{muscle}; changes the size of the pupil\\
\tPupil & the \textbf{hole} that lets light in\\
\tLens & \textbf{convex}; fine-focuses light on the retina; muscles make it fatter or thinner\\
\tRetina & \textbf{light-sensitive cells}; turns light into electrical signals\\
\tOpticNerve & carries the signals to the \textbf{brain}\\
\tBlindSpot & where the optic nerve leaves: \textbf{no} light-sensitive cells\\
\bottomrule
\end{tabularx}
\end{minipage}

\medskip
\textbf{Pupil size:} in \textbf{dim} light the iris makes the pupil \textbf{bigger} to let more light in; in \textbf{bright} light it makes the pupil \textbf{smaller} to protect the retina.

\textbf{Write it in this order} (\emph{How does the eye form an image?})
\begin{enumerate}
\item Light from the object enters through the \textbf{cornea} and the \textbf{pupil}.
\item The cornea and the \textbf{convex lens} \textbf{refract} the light so the rays \textbf{converge}.
\item The rays meet on the \textbf{retina}, forming a \textbf{real}, upside-down image.
\item The retina's cells send electrical signals along the \textbf{optic nerve} to the \textbf{brain}, which turns the image the right way up.
\end{enumerate}

\subsection*{Eye problems and how lenses correct them}
\begin{center}
\FigDefect[5mm]{0}\hspace{3mm}\FigDefect[5mm]{1}\par\medskip
\FigDefect[5mm]{2}\hspace{3mm}\FigDefect[5mm]{3}
\end{center}
\begin{tabularx}{\textwidth}{@{}l X X X@{}}
\toprule
& \textbf{Sees clearly} & \textbf{The problem} & \textbf{Corrected with}\\
\midrule
\tMyopia & \textbf{near} objects; far objects are blurred & rays converge \textbf{too quickly}: focus \textbf{in front of} the retina & \textbf{concave} (diverging) lens\\
\tHyper & \textbf{far} objects; near objects are blurred & rays converge \textbf{too slowly}: focus \textbf{behind} the retina & \textbf{convex} (converging) lens\\
\bottomrule
\end{tabularx}

\section{Lenses}\label{sec:lens}
\begin{minipage}[t]{0.47\textwidth}
\vspace{0pt}
A \textbf{convex} lens is fatter in the middle. It refracts parallel rays so they \textbf{\tConverge} (come together) at the \textbf{\tFocal} F.
The \textbf{\tFocalLength} is the distance from the centre of the lens to F.
The ray through the centre of the lens does not bend.

A \textbf{thicker} (more curved) convex lens bends light \textbf{more}, so its focal point is \textbf{closer} to the lens (shorter focal length).

\textbf{Magnifying glass in sunlight:} rays from the Sun are parallel, so they converge at F; holding the paper at F concentrates the light energy on one small spot.

A \textbf{concave} lens is thinner in the middle. It makes parallel rays \textbf{\tDiverge} (spread out).
\end{minipage}\hfill
\begin{minipage}[t]{0.5\textwidth}
\vspace{0pt}\centering
\FigConvex[5mm]{0}\par\medskip
\FigConvex[5mm]{1}
\end{minipage}

\begin{extbox}[title={Beyond Year 9 \textperiodcentered{} for interest only, not in the Year 9 exam}]\label{sec:beyond}
\small
\textbf{Refractive index.} $n = \dfrac{\text{speed of light in a vacuum}}{\text{speed of light in the material}} = \dfrac{c}{v}$. Glass: $n = \dfrac{300\,000}{200\,000} = 1.5$; water 1.33; diamond 2.42. A larger $n$ means slower light and more bending.

\textbf{Snell's law.} $n_1 \sin\theta_1 = n_2 \sin\theta_2$ ($\theta$ measured from the normal). Going into a material with a \textbf{larger} $n$, $\theta$ gets \textbf{smaller}: the ray bends towards the normal. This is the rule above, written as an equation.

\textbf{Three special rays through a convex lens:}
(1) parallel to the axis $\to$ through F;
(2) through the centre $\to$ straight on;
(3) through F $\to$ parallel to the axis.
Where they meet is the image. A concave lens spreads parallel rays out as if they came from F; a magnifying glass (object inside F) gives a virtual, upright, bigger image.
\begin{center}\FigLensRays[6.2mm]\end{center}

\medskip
\begin{tabularx}{\linewidth}{@{}l l X@{}}
\toprule
\textbf{Object position (convex lens)} & \textbf{Image} & \textbf{Used in}\\
\midrule
beyond 2F & between F and 2F; real, inverted, smaller & eye, camera\\
at 2F & at 2F; real, inverted, same size & photocopier\\
between F and 2F & beyond 2F; real, inverted, bigger & projector\\
at F & no image (rays leave parallel) & spotlight\\
inside F & same side; \textbf{virtual}, upright, bigger & magnifying glass\\
any position (\textbf{concave} lens) & virtual, upright, smaller & glasses for myopia\\
\bottomrule
\end{tabularx}

\begin{center}\FigConcave[5.4mm]\hspace{10mm}\FigMagnifier[5.4mm]\end{center}
\end{extbox}

\section{Explain questions: Achieved, Merit, Excellence}\label{sec:ame}
Explain questions are marked on how far the answer goes. The same question, three levels:

\emph{Explain how a periscope lets you see over a wall.}
\begin{itemize}
\item \textbf{Achieved} (states the idea): ``It uses two mirrors to reflect the light.''
\item \textbf{Merit} (gives the reason): ``Light from the object hits the top mirror and is reflected down the tube, then the bottom mirror reflects it into the eye. The mirrors are at 45\textdegree.''
\item \textbf{Excellence} (links everything): ``Light from the object above the wall hits the top mirror at an \textbf{angle of incidence of 45\textdegree}. By the \textbf{law of reflection} the angle of reflection is also 45\textdegree, so the light turns through \textbf{90\textdegree} and travels down the tube. The bottom mirror is parallel, so the light turns through 90\textdegree{} again and travels horizontally \textbf{into the eye}. The eye receives light from an object it could not see directly.'' (A ray diagram with arrows and the normals supports it.)
\end{itemize}
The chain is always: \textbf{what the light does} $\to$ \textbf{why} (speed, absorbed, reflected, law) $\to$ \textbf{what we see} as a result.

\section{Ways to lose marks}\label{sec:lose}
\begin{enumerate}[label=\textbf{L\arabic*}, leftmargin=9mm]
\item \textbf{\LMtitle{1}.} \LMtext{1}
\item \textbf{\LMtitle{2}.} \LMtext{2}
\item \textbf{\LMtitle{3}.} \LMtext{3}
\item \textbf{\LMtitle{4}.} \LMtext{4}
\item \textbf{\LMtitle{5}.} \LMtext{5}
\item \textbf{\LMtitle{6}.} \LMtext{6}
\item \textbf{\LMtitle{7}.} \LMtext{7}
\item \textbf{\LMtitle{8}.} \LMtext{8}
\end{enumerate}
{\small\emph{These points come from the notes booklet and teaching observation, not from an examiner report.}}

\section{Key words}\label{sec:keywords}
\small
\begin{multicols}{2}
\begin{description}[leftmargin=0pt, labelsep=0.4em, itemsep=0pt]
\item[Light source] emits its own light.
\item[Light reflector] only reflects light.
\item[Medium] a substance light travels through.
\item[Opaque / translucent / transparent] blocks / scatters / transmits light.
\item[Umbra] full shadow; all light blocked.
\item[Penumbra] partial shadow; some light blocked.
\item[Normal] line at 90\textdegree{} to a surface.
\item[Law of reflection] $i = r$.
\item[Plane mirror] a flat mirror.
\item[Real image] can be formed on a screen.
\item[Virtual image] cannot be formed on a screen.
\item[Laterally inverted] left and right swapped.
\item[Refraction] bending of light caused by a change of speed between media.
\item[Optically dense] light travels more slowly in it.
\item[Dispersion] splitting white light into colours.
\item[Spectrum] red, orange, yellow, green, blue, indigo, violet.
\item[Primary colours of light] red, green, blue.
\item[Secondary colours] yellow, cyan, magenta.
\item[Filter] transmits its own colour, absorbs the rest.
\item[Converge / diverge] come together / spread out.
\item[Focal point] where parallel rays meet after a convex lens.
\item[Focal length] lens centre to focal point.
\item[Myopia] short-sighted; concave lens.
\item[Hypermetropia] long-sighted; convex lens.
\end{description}
\end{multicols}

\end{document}
